\chapter{Stargardt Disease}
\index{Stargardt Disease}
\label{sec:Stargardt Disease}

\section{Overview}
\begin{itemize}[noitemsep]
  \item Stargardt disease is the most common inherited macular dystrophy, affecting the central retina and causing progressive central vision loss
  \item It typically begins in childhood or young adulthood
  \item It is a form of juvenile macular degeneration caused by a genetic defect in the ABCA4 gene
\end{itemize}
\section{Causes}
This condition happens because of mutations in the ABCA4 gene, which encodes a protein that clears toxic byproducts (A2E) from photoreceptors. When the protein fails, a toxic substance accumulates in the retinal pigment epithelium, damaging and eventually destroying the photoreceptors in the macula. It is inherited in an autosomal recessive pattern in most cases.
\section{Risk Factors}
\begin{itemize}[noitemsep]
  \item Family history of Stargardt disease (autosomal recessive inheritance)
  \item Having two parents who carry the ABCA4 mutation
  \item Consanguinity increases the chance of inheriting the disease
  \item Sun exposure may accelerate damage
  \item Vitamin A supplementation may worsen the disease
\end{itemize}
\section{Signs and Symptoms}
\subsection{Common symptoms}
\begin{itemize}[noitemsep]
  \item Progressive blurring of central vision in childhood or young adulthood
  \item Difficulty reading, recognizing faces, and seeing fine detail
  \item Central scotomas (blind spots)
  \item Color vision disturbances
  \item Light sensitivity and prolonged adaptation to dim light
  \item Peripheral vision usually remains preserved
\end{itemize}
\subsection{Emergency warning signs}
\begin{itemize}[noitemsep]
  \item Rapidly worsening central vision warrants evaluation to confirm the diagnosis and plan follow-up
\end{itemize}
\section{Progression and Stages}
Vision loss typically begins between ages 6 and 30 and progresses gradually over years to decades. Visual acuity may stabilize at a moderate level of impairment (often 20/200 to 20/400) in adulthood. The disease spares the peripheral retina, so functional peripheral vision and mobility are usually retained.
\section{Complications}
\begin{itemize}[noitemsep]
  \item Severe permanent central vision loss
  \item Difficulty with reading, writing, and other near tasks
  \item Impaired driving and education challenges
  \item Psychological impact and reduced independence
  \item Accelerated progression with supplemental vitamin A
\end{itemize}
\section{Diagnosis}
\begin{itemize}[noitemsep]
  \item Fundus examination showing characteristic flecks and atrophy in the macula (fundus flavimaculatus)
  \item Optical coherence tomography (OCT) showing retinal pigment epithelium changes
  \item Fundus autofluorescence showing abnormal accumulation of lipofuscin
  \item Electroretinography (ERG) to assess photoreceptor function
  \item Genetic testing confirming ABCA4 mutations
\end{itemize}
\section{Prevention}
\begin{itemize}[noitemsep]
  \item Genetic counseling for affected families
  \item Avoiding vitamin A supplementation (which can worsen disease)
  \item Protecting the eyes with UV-blocking sunglasses
  \item Prenatal testing in at-risk families
  \item Participation in registries and research studies
\end{itemize}
\section{Treatment Options}
\begin{itemize}[noitemsep]
  \item There is currently no cure
  \item Gene therapy for ABCA4 mutations is under investigation in clinical trials
  \item Stem cell and retinal transplantation research is ongoing
  \item Low-vision aids and rehabilitation
  \item UV-blocking sunglasses and avoidance of excessive light
  \item Careful dietary counseling regarding vitamin A
\end{itemize}
\section{Living with It}
\begin{itemize}[noitemsep]
  \item Work with low-vision services for magnification aids and training
  \item Use screen-reading and magnification technology
  \item Join support groups and patient communities
  \item Maintain education and employment with appropriate accommodations
  \item Follow research developments in emerging therapies
\end{itemize}
\section{Outlook}
Stargardt disease is progressive, but the rate and final severity vary widely. Most people retain functional peripheral vision and independence, though central vision loss is typically significant. Research into gene and cell therapies offers cautious hope for future treatment.
